% !TEX root = ../main.tex
%% Corrigés — Chapitre 20 — Fonctions de deux variables
%% Résultats vérifiés par verifications/chapitre_20.py.

\corrigesChapitre{20}{Fonctions de deux variables}
\corrigesSouples

\begin{blocLivre}{cor-20-1}
\begin{corrige}{1}{Dérivées partielles : premiers calculs}
\demarche{Dériver par rapport à une variable en gelant l'autre.}
\begin{questions}
  \item $\dfrac{\partial f}{\partial x}=2xy$,
    $\dfrac{\partial f}{\partial y}=x^{2}+3y^{2}$.
  \item $\dfrac{\partial g}{\partial x}=y\,\mathrm{e}^{xy}$,
    $\dfrac{\partial g}{\partial y}=x\,\mathrm{e}^{xy}$.
  \item $\dfrac{\partial h}{\partial x}=\dfrac{2x}{x^{2}+y^{2}}$,
    $\dfrac{\partial h}{\partial y}=\dfrac{2y}{x^{2}+y^{2}}$.
\end{questions}
\end{corrige}
\end{blocLivre}

\begin{blocLivre}{cor-20-2}
\begin{corrige}{2}{Le développement limité d'ordre 1}
\demarche{Le plan tangent remplace la surface, à o(||(h,k)||) près.}
\begin{questions}
  \item $\nabla f(1,2)=(2,4)$ :
    $f(1+h,\,2+k)=5+2h+4k+o\bigl(\norme{(h,k)}\bigr)$.
  \item $f(1{,}02\,;\,1{,}99)\approx5+2(0{,}02)+4(-0{,}01)=\res{5{,}00}$ ;
    la valeur exacte est $5{,}0005$ : erreur de $5\cdot10^{-4}$, l'ordre
    du carré des accroissements.
\end{questions}
\end{corrige}
\end{blocLivre}

\begin{blocLivre}{cor-20-3}
\begin{corrige}{3}{Plans tangents}
\demarche{z - z0 = fx (x - x0) + fy (y - y0) : recopier le gradient.}
\begin{questions}
  \item $\nabla f(1,2)=(2,4)$ :
    $\res{z=5+2(x-1)+4(y-2)}$.
  \item Pour $z=xy$ : $\nabla f(2,3)=(3,2)$ :
    $\res{z=6+3(x-2)+2(y-3)}$.
\end{questions}
\end{corrige}
\end{blocLivre}

\begin{blocLivre}{cor-20-4}
\begin{corrige}{4}{Le gradient et la plus grande pente}
\demarche{Le gradient est la boussole de la montée.}
\begin{questions}
  \item $\nabla f(3,4)=(6,8)$, de norme $\res{10}$.
  \item La croissance est maximale dans la direction du gradient, ici la
    direction \emph{radiale} $(6,8)/10$ : perpendiculaire aux cercles de
    niveau, droit vers l'extérieur — et la pente maximale vaut
    $\norme{\nabla f}=10$.
\end{questions}
\end{corrige}
\end{blocLivre}

\begin{blocLivre}{cor-20-5}
\begin{corrige}{5}{Lignes de niveau}
\demarche{Résoudre f = c pour chaque constante : la carte se dessine.}
Pour $f$ : les cercles $x^{2}+y^{2}=c$ ($c>0$), réduits à l'origine pour
$c=0$. Pour $g$ : les hyperboles $xy=c$, et la réunion des deux axes pour
$c=0$. Pour $h$ : les hyperboles $x^{2}-y^{2}=c$, et les deux
bissectrices pour $c=0$. Par $(1,2)$ passent le cercle de rayon
$\sqrt5$, l'hyperbole $xy=2$ et l'hyperbole $x^{2}-y^{2}=-3$.
\end{corrige}
\end{blocLivre}

\begin{blocLivre}{cor-20-6}
\begin{corrige}{6}{Dérivée selon un vecteur}
\demarche{Un produit scalaire avec le gradient — chapitre 18 en action.}
\begin{questions}
  \item $\nabla f(1,1)=(2,1)$ et
    $D_uf(1,1)=\dfrac{3\times2+4\times1}{5}=\res{2}$.
  \item Par Cauchy-Schwarz,
    $D_uf=\langle\nabla f,u\rangle\leqslant\norme{\nabla f}=\sqrt5$,
    avec égalité pour $u=\dfrac{(2,1)}{\sqrt5}$ : la plus grande pente,
    de valeur $\res{\sqrt5}$.
\end{questions}
\end{corrige}
\end{blocLivre}

\begin{blocLivre}{cor-20-7}
\begin{corrige}{7}{La règle de la chaîne}
\demarche{Gradient le long du chemin, produit scalaire avec la vitesse.}
\begin{questions}
  \item $\bigl(f\circ\gamma\bigr)'(t)
    =2\cos t\,(-\sin t)+2\sin t\,\cos t=\res{0}$ : prévisible, le cercle
    unité est une \emph{ligne de niveau} de $f$ ($f\circ\gamma\equiv1$).
  \item Direct : $(\cos t\sin t)'=\cos2t$. Par la chaîne :
    $y\,(-\sin t)+x\,\cos t=\cos^{2}t-\sin^{2}t=\cos2t$ : accord
    parfait.
\end{questions}
\end{corrige}
\end{blocLivre}

\begin{blocLivre}{cor-20-8}
\begin{corrige}{8}{Premiers points critiques}
\demarche{Gradient nul d'abord, étude locale ensuite — le col guette.}
Les trois fonctions ont $(0,0)$ pour unique point critique.
$x^{2}+y^{2}$ : \res{\text{minimum global}} (valeurs positives).
$x^{2}-y^{2}$ : $f(t,0)=t^{2}>0>-t^{2}=f(0,t)$ :
\res{\text{point-selle}}, pas d'extremum. $xy$ : $f(t,t)=t^{2}>0>
f(t,-t)$ : point-selle également — le gradient nul ne promet rien.
\end{corrige}
\end{blocLivre}

\begin{blocLivre}{cor-20-9}
\begin{corrige}{9}{Dérivées partielles sans continuité}
\demarche{Deux directions ne suffisent pas : la diagonale raconte une
autre histoire que les axes.}
\begin{questions}
  \item Le long de $y=0$ : $f(x,0)=0\to0$. Le long de $y=x$ :
    $f(x,x)=\dfrac{x^{2}}{2x^{2}}=\dfrac12$ : deux limites différentes,
    $f$ n'est \res{\text{pas continue}} en $(0,0)$.
  \item $\dfrac{f(h,0)-f(0,0)}{h}=0\to0$ : les deux dérivées partielles
    existent et valent $0$. Il manque la \emph{continuité des dérivées
    partielles} (classe $\mathcal{C}^{1}$) : seule elle garantit le
    développement limité — les dérivées partielles seules ne contrôlent
    que deux directions sur une infinité.
\end{questions}
\end{corrige}
\end{blocLivre}

\begin{blocLivre}{cor-20-10}
\begin{corrige}{10}{Composées et dérivées partielles}
\demarche{La chaîne, variable par variable : chaque argument transmet sa
dérivée.}
\begin{questions}
  \item $\dfrac{\partial F}{\partial u}
    =\dfrac{\partial f}{\partial x}(u+v,\,uv)
    +v\,\dfrac{\partial f}{\partial y}(u+v,\,uv)$, et symétriquement
    $\dfrac{\partial F}{\partial v}
    =\dfrac{\partial f}{\partial x}+u\,\dfrac{\partial f}{\partial y}$.
  \item Avec $f(x,y)=x^{2}y+y^{3}$ : le développement complet de
    $F(u,v)=(u+v)^{2}uv+(uv)^{3}$ et la formule donnent le même
    polynôme — vérification faite, la règle est fiable.
\end{questions}
\end{corrige}
\end{blocLivre}

\begin{blocLivre}{cor-20-11}
\begin{corrige}{11}{Le gradient orthogonal aux lignes de niveau}
\demarche{Constante le long du chemin : dérivée nulle, produit scalaire
nul.}
\begin{questions}
  \item $f\bigl(\gamma(t)\bigr)=3\cos^{2}t+2\cdot\dfrac32\sin^{2}t=3$ :
    l'ellipse est la ligne de niveau $3$.
  \item $\nabla f(\gamma)=\bigl(2\sqrt3\cos t,\,4\sqrt{3/2}\sin t\bigr)$
    et $\gamma'=\bigl(-\sqrt3\sin t,\,\sqrt{3/2}\cos t\bigr)$ : le
    produit scalaire vaut $-2\cdot3\ldots+2\cdot3\ldots=0$ — mais sans
    calcul : $f\circ\gamma$ est constante, donc
    $\langle\nabla f(\gamma),\gamma'\rangle=(f\circ\gamma)'=0$ : le
    gradient est \res{\text{orthogonal aux lignes de niveau}}.
\end{questions}
\end{corrige}
\end{blocLivre}

\begin{blocLivre}{cor-20-12}
\begin{corrige}{12}{Dériver le long d'un arc}
\demarche{Deux routes, un résultat : c'est la règle de la chaîne qui se
vérifie.}
\begin{questions}
  \item $f\bigl(\gamma(t)\bigr)=t^{2}\cdot t^{2}=t^{4}$ : dérivée
    $\res{4t^{3}}$.
  \item $\nabla f(t,t^{2})=(2t^{3},\,t^{2})$ et
    $\gamma'(t)=(1,2t)$ : produit scalaire
    $2t^{3}+2t^{3}=4t^{3}$ — identique.
\end{questions}
\end{corrige}
\end{blocLivre}

\begin{blocLivre}{cor-20-13}
\begin{corrige}{13}{Des équations aux dérivées partielles}
\demarche{Dérivée partielle nulle : constante par rapport à cette variable,
fonction arbitraire de l'autre.}
\begin{questions}
  \item À $y$ fixé, $x\mapsto f(x,y)$ a une dérivée nulle sur $\R$ :
    elle est constante, et $\res{f(x,y)=\varphi(y)}$ avec $\varphi$ de
    classe $\mathcal{C}^{1}$ arbitraire.
  \item $\dfrac{\partial f}{\partial x}=y$ donne
    $f(x,y)=xy+\varphi(y)$ ; la seconde équation impose
    $x+\varphi'(y)=x$, soit $\varphi$ constante :
    $\res{f(x,y)=xy+c}$ — à constante près, la fonction est
    reconstituée de ses dérivées.
\end{questions}
\end{corrige}
\end{blocLivre}

\begin{blocLivre}{cor-20-14}
\begin{corrige}{14}{Le col et le sommet}
\demarche{Résoudre le gradient, puis faire parler des chemins bien
choisis.}
\begin{questions}
  \item $\nabla f=0$ s'écrit $y=x^{2}$ et $x=y^{2}$ : en substituant,
    $x=x^{4}$, d'où les points critiques $\res{(0,0)}$ et
    $\res{(1,1)}$.
  \item En $(0,0)$ : $f(t,0)=t^{3}$ change de signe : pas d'extremum
    (un col). En $(1,1)$, posons $x=1+h$, $y=1+k$ :
    \[
    f(1+h,1+k)+1=3\bigl(h^{2}-hk+k^{2}\bigr)+h^{3}+k^{3},
    \]
    et $h^{2}-hk+k^{2}=\bigl(h-\tfrac k2\bigr)^{2}+\tfrac34k^{2}>0$ hors
    de l'origine : le terme quadratique domine localement,
    $(1,1)$ est un \res{\text{minimum local, de valeur }-1}.
\end{questions}
\end{corrige}
\end{blocLivre}

\begin{blocLivre}{cor-20-15}
\begin{corrige}{15}{Extremums sur un carré}
\demarche{L'intérieur par le gradient, le bord segment par segment.}
\begin{questions}
  \item $\nabla f=(2x-1,\,2y-1)=0$ :
    $\Bigl(\dfrac12,\dfrac12\Bigr)$, où
    $f=-\dfrac12$.
  \item Sur $y=0$ : $x^{2}-x$, minimum $-\frac14$ en $\frac12$, maximum
    $0$ aux extrémités — et de même sur les trois autres côtés (symétrie
    $x\leftrightarrow y$, $t\leftrightarrow1-t$). Bilan :
    \res{\text{minimum global }-\tfrac12} au centre,
    \res{\text{maximum global }0} atteint aux \emph{quatre coins} du
    carré.
\end{questions}
\end{corrige}
\end{blocLivre}

\begin{blocLivre}{cor-20-16}
\begin{corrige}{16}{Le plan tangent comme calculette}
\demarche{Le DL d'ordre 1 est la meilleure calculatrice de poche.}
\begin{questions}
  \item $\nabla f(3,4)=\Bigl(\dfrac35,\dfrac45\Bigr)$ — un vecteur
    unitaire : la norme croît à vitesse $1$ le long des rayons.
  \item $\sqrt{3{,}03^{2}+3{,}96^{2}}\approx5+\dfrac35(0{,}03)
    +\dfrac45(-0{,}04)=\res{4{,}986}$ ; la valeur exacte est
    $4{,}98623$ : quatre décimales sûres.
\end{questions}
\end{corrige}
\end{blocLivre}

\begin{blocLivre}{cor-20-17}
\begin{corrige}{17}{L'identité d'Euler}
\demarche{Dériver l'homogénéité par rapport au facteur d'échelle.}
\begin{questions}
  \item $(\lambda x)^{2}(\lambda y)=\lambda^{3}x^{2}y$ et
    $(\lambda x)^{3}-2(\lambda x)(\lambda y)^{2}
    =\lambda^{3}\bigl(x^{3}-2xy^{2}\bigr)$ : degré $3$.
  \item Soit $g(t)=f(tx,ty)=t^{d}f(x,y)$. La règle de la chaîne donne
    $g'(t)=x\dfrac{\partial f}{\partial x}(tx,ty)
    +y\dfrac{\partial f}{\partial y}(tx,ty)$, et la dérivée directe
    $g'(t)=d\,t^{d-1}f(x,y)$ ; en $t=1$ :
    \[
    \res{x\frac{\partial f}{\partial x}+y\frac{\partial f}{\partial y}
    =d\,f} .
    \]
\end{questions}
\end{corrige}
\end{blocLivre}

\begin{blocLivre}{cor-20-18}
\begin{corrige}{18}{La droite des moindres carrés}
\demarche{Annuler le gradient d'une somme de carrés : deux équations
linéaires.}
\begin{questions}
  \item $\dfrac{\partial E}{\partial a}=0$ et
    $\dfrac{\partial E}{\partial b}=0$ s'écrivent
    $5a+3b=3$ et $3a+3b=2$ :
    $\res{a=\dfrac12,\quad b=\dfrac16}$.
  \item La droite $y=\dfrac x2+\dfrac16$, avec l'erreur minimale
    $E=\dfrac16$. C'est bien un minimum : $E$ est une fonction
    quadratique positive qui tend vers $+\infty$ à l'infini, et ne
    possède qu'un point critique — la droite de régression des
    statisticiens.
\end{questions}
\end{corrige}
\end{blocLivre}

\begin{blocLivre}{cor-20-19}
\begin{corrige}{19}{Le théorème de Schwarz, et sa limite}
\demarche{Pour les fonctions régulières, l'ordre de dérivation est
indifférent — pour les autres, non.}
\begin{questions}
  \item Les deux dérivées croisées de
    $\mathrm{e}^{xy}\sin(x+y)$ coïncident (calcul direct) : c'est le
    théorème de Schwarz, valable pour toute fonction de classe
    $\mathcal{C}^{2}$.
  \item Pour $g$, un calcul (ou une vérification numérique) donne en
    $(0,0)$ :
    $\dfrac{\partial^{2}g}{\partial y\partial x}=-1$ et
    $\dfrac{\partial^{2}g}{\partial x\partial y}=+1$ : sans la
    continuité des dérivées secondes, \res{\text{l'ordre de dérivation
    compte}} — la régularité n'est pas une politesse de rédaction.
\end{questions}
\end{corrige}
\end{blocLivre}

\begin{blocLivre}{cor-20-20}
\begin{corrige}{20}{Fonctions harmoniques}
\demarche{Deux dérivées secondes, une somme nulle — et les complexes en
coulisse.}
\begin{questions}
  \item $\Delta(x^{2}-y^{2})=2-2=0$ ; $\Delta(xy)=0$ ;
    $\Delta\bigl(\mathrm{e}^{x}\cos y\bigr)
    =\mathrm{e}^{x}\cos y-\mathrm{e}^{x}\cos y=0$ ; et pour
    $\ln\bigl(x^{2}+y^{2}\bigr)$, le calcul donne également $0$ hors de
    l'origine : quatre fonctions \res{\text{harmoniques}}.
  \item $x^{2}-y^{2}$ et $2xy$ sont les parties réelle et imaginaire de
    $z^{2}=(x+\mathrm{i}y)^{2}$ (chapitre 3) — comme
    $\mathrm{e}^{x}\cos y=\operatorname{Re}\bigl(\mathrm{e}^{z}\bigr)$ :
    les fonctions de la variable complexe fabriquent des harmoniques en
    série, porte ouverte sur l'analyse complexe de deuxième année.
\end{questions}
\end{corrige}
\end{blocLivre}

\begin{blocLivre}{cor-20-21}
\begin{corrige}{21}{Une étude complète}
\demarche{Pour le global : minorer par deux problèmes à une variable.}
\begin{questions}
  \item $\nabla f=\bigl(4x^{3}-4(x-y),\,4y^{3}+4(x-y)\bigr)$ ; en
    sommant, $x^{3}+y^{3}=0$, donc $y=-x$, puis $4x^{3}-8x=0$ : points
    critiques $\res{(0,0)}$ et $\res{\pm(\sqrt2,-\sqrt2)}$.
  \item En $(0,0)$ : $f(t,t)=2t^{4}>0$ et
    $f(t,-t)=2t^{4}-8t^{2}<0$ près de $0$ : col. Pour le global :
    $(x-y)^{2}\leqslant2\bigl(x^{2}+y^{2}\bigr)$ donne
    \[
    f(x,y)\geqslant\bigl(x^{4}-4x^{2}\bigr)+\bigl(y^{4}-4y^{2}\bigr)
    \geqslant-4-4=-8 ,
    \]
    car $t^{4}-4t^{2}\geqslant-4$. Et $f\bigl(\sqrt2,-\sqrt2\bigr)=-8$ :
    le \res{\text{minimum global vaut }-8}, atteint exactement en
    $\pm(\sqrt2,-\sqrt2)$.
\end{questions}
\end{corrige}
\end{blocLivre}

\begin{blocLivre}{cor-20-22}
\begin{corrige}{22}{Distance d'un point à une parabole}
\demarche{Une variable suffit — et le gradient explique la géométrie.}
\begin{questions}
  \item $d(x)=x^{2}+\bigl(x^{2}-1\bigr)^{2}$ a pour dérivée
    $2x\bigl(2x^{2}-1\bigr)$ : candidats $0$ et
    $\pm\dfrac{1}{\sqrt2}$, avec $d(0)=1$ et
    $d\bigl(\pm\tfrac{1}{\sqrt2}\bigr)=\dfrac34$ : la distance vaut
    $\res{\dfrac{\sqrt3}{2}}$.
  \item Au point $\bigl(\tfrac{1}{\sqrt2},\tfrac12\bigr)$, le vecteur
    joignant $(0,1)$ au point est
    $\bigl(\tfrac{1}{\sqrt2},-\tfrac12\bigr)$ et la tangente a pour
    direction $(1,\sqrt2)$ : produit scalaire
    $\tfrac{1}{\sqrt2}-\tfrac{\sqrt2}{2}=0$ — le plus court chemin
    arrive \emph{perpendiculairement} à la courbe, comme la projection
    orthogonale du chapitre 18.
\end{questions}
\end{corrige}
\end{blocLivre}

\begin{blocLivre}{cor-20-23}
\begin{corrige}{23}{Passage en coordonnées polaires}
\demarche{La chaîne appliquée au changement de variables le plus utile du
plan.}
\begin{questions}
  \item $\dfrac{\partial g}{\partial r}
    =\cos\theta\,\dfrac{\partial f}{\partial x}
    +\sin\theta\,\dfrac{\partial f}{\partial y}$ et
    $\dfrac{\partial g}{\partial\theta}
    =-r\sin\theta\,\dfrac{\partial f}{\partial x}
    +r\cos\theta\,\dfrac{\partial f}{\partial y}$.
  \item Pour $f=x^{2}+y^{2}$ : $g=r^{2}$, donc
    $\dfrac{\partial g}{\partial r}=2r$ et
    $\dfrac{\partial g}{\partial\theta}=0$ — prévisible : $f$ est
    \emph{radiale}, constante sur les cercles ; les polaires épousent
    ses lignes de niveau.
\end{questions}
\end{corrige}
\end{blocLivre}

\begin{blocLivre}{cor-20-24}
\begin{corrige}{24}{Moindres carrés, version matricielle}
\demarche{Gradient, équations normales, projection : trois chapitres pour
un même minimum.}
\begin{questions}
  \item $\norme{AX-B}^{2}=(AX-B)^{\mathsf T}(AX-B)$ a pour gradient
    $2A^{\mathsf T}(AX-B)$ : il s'annule exactement pour
    $A^{\mathsf T}AX=A^{\mathsf T}B$, soit
    $\begin{pmatrix}5&3\\3&3\end{pmatrix}X
    =\begin{pmatrix}3\\2\end{pmatrix}$ :
    $\res{X=\Bigl(\dfrac12,\ \dfrac16\Bigr)}$ — la droite de
    l'exercice 18, retrouvée matriciellement.
  \item Le résidu $AX-B$ vérifie $A^{\mathsf T}(AX-B)=0$ : il est
    orthogonal aux colonnes de $A$, donc $AX$ est la
    \res{\text{projection orthogonale}} de $B$ sur
    $\operatorname{Im}A$ (chapitre 18), et le coût minimal vaut
    $\dfrac16$.
  \item Le même problème s'est laissé résoudre par le gradient
    (chapitre 20), par les équations normales (chapitre 13) et par la
    projection orthogonale (chapitre 18) : trois langages, une seule
    géométrie — c'est sur cette convergence que s'achève le cours.
\end{questions}
\end{corrige}
\end{blocLivre}

% !TEX root = ../../main.tex
%% Corrigé du défi du chapitre 20 — vérifié par verifications/defis_16_20.py.

\begin{blocLivre}{cor-20-25a}
\begin{corrige}{25}{Le point de Fermat}
\demarche{Le gradient d'une distance est un vecteur unitaire : annuler
$\nabla f$, c'est équilibrer trois vecteurs unitaires — d'où les angles de
$120^{\circ}$. Une dérivée directionnelle négative écarte les sommets.}
\partieExo{A}{Gradient et points critiques}
\begin{questions}
  \item Avec $P=(a,b)$ : $d_P(x,y)=\sqrt{(x-a)^{2}+(y-b)^{2}}$, racine
    d'une fonction polynomiale strictement positive sur
    $\R^{2}\setminus\ens{P}$ : $d_P$ y est $\mathcal{C}^{1}$, avec
    $\frac{\partial d_P}{\partial x}=\frac{x-a}{d_P}$ et
    $\frac{\partial d_P}{\partial y}=\frac{y-b}{d_P}$, soit
    $\res{\nabla d_P(M)=\frac{M-P}{\norme{M-P}}}$, de norme $1$.
  \item $f=d_A+d_B+d_C$ est $\mathcal{C}^{1}$ sur $U$ et
    $\nabla f(M)=u_A+u_B+u_C$ : les points critiques sont ceux où cette
    somme de vecteurs unitaires est \res{\text{nulle}}.
  \item Si $u+v+w=0$ avec $\norme u=\norme v=\norme w=1$ :
    $1=\norme{u+v}^{2}=2+2\langle u,v\rangle$, d'où $\langle u,v\rangle=-\frac12$
    et un angle de $\frac{2\pi}{3}$ ; de même pour les autres couples.
    L'angle $\widehat{AMB}$ est celui de $-u_A$ et $-u_B$, donc celui de
    $u_A$ et $u_B$ : \res{\tfrac{2\pi}{3}}, comme les deux autres.
\end{questions}
\partieExo{B}{Existence et position du minimum}
\begin{questions}[resume]
  \item $MP\geqslant\norme M-\norme P$, d'où la minoration. Si
    $\norme M>R=\frac13\bigl(\norme A+\norme B+\norme C+f(O)\bigr)$, alors
    $f(M)>f(O)$. Sur le disque fermé de centre $O$ et de rayon $R$, $f$
    continue atteint un minimum $f(M_0)\leqslant f(O)$ ; hors du disque,
    $f(M)>f(O)\geqslant f(M_0)$ : $f(M_0)$ est un
    \res{\text{minimum global}}.
  \item $d_B$ et $d_C$ sont dérivables en $A$, et
    $\frac{\mathrm{d}}{\mathrm{d}t}f(A+tu)\big|_{t=0^{+}}
    =1+\langle u,\nabla d_B(A)\rangle+\langle u,\nabla d_C(A)\rangle
    =1-\langle u,e_B+e_C\rangle$, où $e_B$, $e_C$ sont les vecteurs
    unitaires de $A$ vers $B$ et $C$. Comme
    $\norme{e_B+e_C}^{2}=2+2\cos\widehat A=4\cos^{2}\frac{\widehat A}{2}$, le
    choix $u=\frac{e_B+e_C}{\norme{e_B+e_C}}$ donne la dérivée
    $1-2\cos\frac{\widehat A}{2}$, strictement négative car
    $\frac{\widehat A}{2}<\frac\pi3$ : $f$ décroît en partant de $A$, qui
    n'est pas un minimum. Le minimum est donc atteint dans l'ouvert $U$,
    en un \res{\text{point critique}} $F$ : les côtés y sont vus sous des
    angles de $\frac{2\pi}{3}$.
\end{questions}
\end{corrige}
\end{blocLivre}

\begin{blocLivre}{cor-20-25b}
\begin{corrigeSuite}
\partieExo{C}{Un exemple, et le retour de Napoléon}
\begin{questions}[start=6]
  \item $\cos\widehat A=\frac{\langle(2,0),(1,2)\rangle}{2\sqrt5}=\frac{1}{\sqrt5}$,
    soit $\widehat A=\widehat B\approx63{,}4^{\circ}$ et
    $\widehat C\approx53{,}1^{\circ}$ : tous inférieurs à $120^{\circ}$.
    Pour $F=(0,y)$, $\widehat{AFB}=\frac{2\pi}{3}$ équivaut à
    $\tan\frac\pi3=\frac1y$, soit $y=\frac{1}{\sqrt3}$ ; alors $u_A$, $u_B$,
    $u_C$ ont pour angles polaires $30^{\circ}$, $150^{\circ}$, $270^{\circ}$ :
    leur somme est nulle, $F$ est bien le point critique, et
    \[
    f(F)=2\sqrt{1+\tfrac13}+\Bigl(2-\tfrac{1}{\sqrt3}\Bigr)
    =\tfrac{4}{\sqrt3}-\tfrac{1}{\sqrt3}+2=\res{2+\sqrt3}\approx3{,}732 .
    \]
  \item Le triangle équilatéral de côté $AB=2$ a pour hauteur $\sqrt3$ :
    $C'=(0,-\sqrt3)$. Les points $C'$, $F$, $C$ sont sur l'axe, dans cet
    ordre : $F\in[CC']$, et $CC'=2+\sqrt3=\res{f(F)}$. On retrouve le défi
    du chapitre 3 : les droites $(AA')$, $(BB')$, $(CC')$ concourent au point
    de Fermat, et leur longueur commune est la longueur minimale du
    réseau.
  \item Le gradient de $f$ est orthogonal aux lignes de niveau (chapitre 20)
    et s'annule en $F$, où les lignes se resserrent en petites courbes
    fermées autour du minimum. Loin du triangle,
    $f(M)\approx3\,\norme{M-G}$ ($G$ centre de gravité) : les lignes de
    niveau ressemblent à des cercles centrés près de $G$.
\end{questions}
\end{corrigeSuite}
\end{blocLivre}

